% !PRINTED 116
\weberpage{116}
%<<CONT

\[
 y=\pm\sqrt{-a},
\]
und folglich
\[
 x=-\frac{a_1}{2}\pm\sqrt{\frac{a_1^2}{4}-a_2}
   =\frac{-a_1\pm\sqrt{a_1^2-4a_2}}{2}.
\]
Wir haben also zwei Wurzeln
\begin{equation}
\alpha=\frac{-a_1+\sqrt{a_1^2-4a_2}}{2},\qquad
\beta =\frac{-a_1-\sqrt{a_1^2-4a_2}}{2}.
\tag{5}
\end{equation}

\begin{center}
{\large \S.\ 35.}

\vspace{0.5em}
{\large\gesperrt{Cubische Gleichungen. Cardanische Formel.}}
\end{center}

Die Gleichung dritten Grades nehmen wir in der reducirten Form an
\begin{equation}
y^3+ay+b=0
\tag{1}
\end{equation}
und führen zwei neue Unbekannte $u,v$ ein, indem wir
\begin{equation}
y=u+v
\tag{2}
\end{equation}
setzen. Dies in (1) eingesetzt, ergiebt
\begin{equation}
u^3+v^3+(3uv+a)(u+v)+b=0.
\tag{3}
\end{equation}

Wir bestimmen eine der beiden Grössen $u,v$ durch die andere nach der Gleichung
\begin{equation}
3uv=-a,
\tag{4}
\end{equation}
und erhalten aus (3)
\begin{equation}
u^3+v^3=-b.
\tag{5}
\end{equation}
Aus (4) und (5) aber lassen sich $u^3$ und $v^3$ durch eine Quadratwurzel bestimmen. Man erhält nämlich
\begin{equation}
(u^3-v^3)^2=(u^3+v^3)^2-4u^3v^3
=b^2+\frac{4a^3}{27}.
\tag{6}
\end{equation}
Setzen wir zur Abkürzung
\begin{equation}
R=\frac{b^2}{4}+\frac{a^3}{27},
\tag{7}
\end{equation}
so findet sich
\[
 u^3-v^3=2\sqrt R.
\]
